% SPDX-License-Identifier: Apache-2.0
% covers: ExMon
\datasheet{ExMon}{The exclusive monitor: reservations for two ports, and the answer to each exclusive write.}

\dsfacts{\code{lib/parts/src/bus/exmon.rs}}%
{\code{txhdl\_parts::bus::exmon::ExMon<J, I, P0, P1, RB, RM>}}%
{\code{//docs:axi}}

\subsection*{Function}

The monitor makes \code{lr.w}, \code{sc.w} and the AMOs of two harts
atomic against each other and against every other host on the memory
(issue 1408). It stands on an AXI4 link after the arbiter, where every
write passes in the order the peripherals take it, and watches two of
the arbiter's ports, \code{P0} and \code{P1}.

An exclusive read, one with \code{lock} set, from a watched port makes
a reservation for that port of the sixteen-byte line it reads. A write
into the range that passes clears every reservation whose line its
burst covers, whoever made it. An exclusive write is the other half.
It passes when its port's reservation stands for its line and it is
one beat, and its answer goes back as \code{ExOkay}, which is how the
host learns that it stored. Otherwise it fails: the monitor takes its
beats, writes nothing, and answers \code{Okay} itself, so the write
never reaches the peripheral or anything between. An exclusive write,
whether it stores or not, uses its port's reservation up, as \code{sc.w} does.

The monitor keeps exclusives for one range of addresses, the memory's:
an address \code{a} with \code{(a \& RM) == RB}. An exclusive read
outside it makes no reservation, and an exclusive write outside it
fails.

It passes \code{lock} on as zero, since the peripheral behind it does
not keep exclusives itself. It has no read data channel: the read data
goes back around it.

\subsection*{Parameters}

\begin{center}\footnotesize
\begin{tabular}{@{}lp{0.7\textwidth}@{}}
\toprule
Parameter & Meaning \\
\midrule
\code{J} & The identifier width on the link, in bits, at most five:
\code{wids} keeps a bit for each identifier, 32 of them. \\
\code{I} & The host's own bits at the bottom of the identifier; the
port is the identifier shifted right by \code{I}. \\
\code{P0} & The first watched port. \\
\code{P1} & The second watched port. \\
\code{RB} & The base of the range kept. \\
\code{RM} & The mask of the range kept. \\
\bottomrule
\end{tabular}
\end{center}

Addresses and data are 32 bits, and the write strobe four bits, fixed
in the source. The tables below are generated for the board's,
\code{ExMon<5, 2, 0, 7, 0x4000\_0000, 0xc000\_0000>}: identifiers of
five bits with two of the host's own, the harts on ports 0 and 7, and
the range of the DDR3.

\dstables{ExMon}

\subsection*{Behaviour}

\begin{itemize}
\item \textbf{Reads.} A read passes when the link onward has room. An
exclusive read into the range also waits until every write into the
range the monitor has passed is answered, which \code{wids} records by
identifier. So it reads what those writes wrote, and no write in flight
when it was taken lands after it. Writes outside the range, which a slow
peripheral may hold for long, it does not wait for.
\item \textbf{Reservations.} An exclusive read into the range from port
\code{P0} sets \code{v0} and keeps the line, the address above its four
low bits, in \code{line0}; from \code{P1}, \code{v1} and \code{line1}.
From any other port it makes none.
\item \textbf{Clears.} A write into the range that passes is kept in
\code{fwv}, \code{fwline} and \code{fwcnt}, its first line and how many
it covers. The cycle after, every reservation in those lines is
cleared. That clear wins over a reservation set in the write's own
cycle.
\item \textbf{Write address.} One burst is taken at a time
(\code{wpend}), and none while a failed write's answer waits
(\code{binj}). A plain write passes when the link onward has room.
\item \textbf{The decision.} An exclusive write waits a cycle at the
head, where it is decided into \code{xk\_v}, \code{xk\_ok} and
\code{xk\_p0}, and goes on the next, so only registers stand before the
router. It stores when its port is \code{P0} or \code{P1}, that port's
reservation stands, is not being cleared, and is for its line, it is
one beat (\code{len} zero), and it is in the range.
\item \textbf{A stored exclusive write} passes, and \code{x0} or
\code{x1} keeps its identifier. Its answer, if \code{Okay}, goes back
as \code{ExOkay}. As a write into its own line it clears its port's
reservation the cycle after.
\item \textbf{A failed exclusive write} is taken and not passed on;
\code{wdrop} marks its beats, which are taken and dropped. It clears its
port's reservation and adds one to \code{fails}. After its last beat
\code{binj} holds its answer, \code{Okay} with its identifier in
\code{binjid}. That answer is given once every write into the range is
answered, so it never overtakes an answer to a write its host sent
before it.
\item \textbf{Answers.} The peripheral's answers pass otherwise
unchanged, and each clears its identifier's bit in \code{wids}. A failed
write's answer goes first when both are ready.
\item \textbf{\code{xdup}} is set for good when a write passes under the
identifier of an exclusive write whose answer is still to come, which
would make that answer ambiguous. A test reads it.
\item \textbf{\code{alternate}} is a test's: set, it fails every other
exclusive write, the first included, as if its reservation were gone,
with \code{flip} keeping which (issue 1474). Nothing in the design sets
it.
\end{itemize}

\subsection*{Verification}

\begin{itemize}
\item Six tests in \code{bus::exmon}, in \code{//lib/parts:parts\_test},
run two hosts, an arbiter and the monitor before a memory that takes a
read's words at its address. Every run checks that \code{xdup} stays
clear.
\begin{itemize}
\item \code{a\_pair\_with\_nothing\_between\_stores} gets \code{ExOkay}
and the word stored.
\item \code{a\_write\_between\_breaks\_the\_pair}, with the arbiter's
hold not built, gets \code{Okay} and nothing stored.
\item \code{an\_exclusive\_read\_waits\_for\_a\_write\_in\_flight} reads
the word of another host's burst still being written, and the pair
keeps.
\item \code{a\_failed\_writes\_answer\_overtakes\_nothing} checks a
failed write's answer comes after the host's earlier write's.
\item \code{a\_failed\_exclusive\_write\_uses\_the\_reservation\_up}
checks that after a failed \code{sc} of one word an \code{sc} of the
reserved word fails too.
\item \code{alternate\_fails\_every\_other\_exclusive\_write} checks the
first try fails and the second stores.
\end{itemize}
\item \code{//cpu/vreteno:board\_test} runs the board.
\begin{itemize}
\item Three read \code{fails}:\\
\code{the\_a\_instructions\_on\_the\_ddr3\_are\_exclusive\_pairs},\\
\code{another\_hosts\_write\_breaks\_a\_reservation} and\\
\code{amos\_lose\_no\_count\_to\_the\_other\_hosts\_writes}.
\item Two set \code{alternate}:\\
\code{every\_amo\_reads\_again\_ending\_a\_line} and\\
\code{an\_amo\_reads\_again\_while\_a\_walk\_is\_due}.
\item \code{two\_harts\_lose\_no\_count} has both harts count with
\code{amoadd.w} and with \code{lr.w} and \code{sc.w}.
\end{itemize}
\item \code{//cpu/vreteno/act:vreteno\_excl\_tests} runs the A
extension's architectural tests with the monitor between the hart's
tracker and the memory, which \code{act\_run} sets up for those
tests.
\end{itemize}

\subsection*{Used by}

\code{Board}, as \code{exmon}, an
\code{ExMon<5, 2, 0, 7, 0x4000\_0000, 0xc000\_0000>} between the arbiter
and the router: the harts are on the arbiter's ports 0 and 7, and the
range is the DDR3's. \code{act\_run} uses
\code{ExMon<2, 2, 0, 1, 0x4000\_0000, 0xc000\_0000>} with one host.

\subsection*{Limits}

Two ports are watched; an exclusive read from any other makes no
reservation. One range is kept. A reservation is a line of sixteen
bytes, and an exclusive write of more than one beat fails. Write bursts
pass one at a time. \code{J} is at most five, and nothing checks it.
The monitor relies on the arbiter's hold to keep new writes from other
hosts out while a pair is open, so that the wait of an exclusive read
ends.
